
\documentclass[pdflatex,sn-vancouver-num]{sn-jnl}
\usepackage{booktabs}
\usepackage{amsmath}
\begin{document}
\title[Article Title]{BMC Bioinformatics: Official Springer Nature Submission Preview}
\author*[1]{\fnm{A.} \sur{Placeholder}}\email{verified@example.org}
\author[2]{\fnm{B.} \sur{Example}}
\affil*[1]{\orgdiv{Department}, \orgname{Placeholder Institute}, \orgaddress{\city{Placeholder City}, \country{Country}}}
\affil[2]{\orgdiv{Laboratory}, \orgname{Example University}, \orgaddress{\city{Example City}, \country{Country}}}
\abstract{This submission-style preview is rendered with the imported official Springer Nature sn-article template package. The content remains placeholder text, but the page is intentionally structured as a realistic article so front matter, section hierarchy, figure-table rhythm, and bibliography style are easier to inspect.}
\keywords{Springer Nature template, official preview, submission layout, figure placement, table placement}
\maketitle
\section{Introduction}
The introduction is written as compact technical prose so the preview reads more like a practical Springer Nature submission than a bare template check page.
\section{Related Work}
The related-work section is explicit so the preview reveals how literature framing, methods, and experiments separate within the article.
\section{Methodology}
The methodology section introduces representative mathematics and a figure placeholder so the imported sn-jnl package can be inspected under more realistic technical content.
\begin{equation}
\mathcal{L} = \sum_{i=1}^{M} w_i \left\| y_i - \hat{y}_i \right\|_2^2.
\end{equation}
\begin{figure}[t]
\centering
\fbox{\parbox[c][2.0in][c]{0.85\linewidth}{\centering Springer Figure Placeholder\\[0.4em]\normalsize Pipeline overview, experiment setup, or system diagram}}
\caption{Representative figure block used to inspect the official Springer Nature template, caption spacing, and the transition from methods text to the first major figure.}
\end{figure}
\section{Experiments}
The experiments section provides enough structure to expose table treatment, subsection rhythm, and the transition from methods to quantitative comparison.
\subsection{Experimental Setup}
This paragraph acts as a placeholder for datasets, experimental conditions, baseline systems, and protocol details.
\subsection{Quantitative Results}
\begin{table}[t]
\caption{Quantitative comparison block used to inspect the official Springer Nature template under realistic submission-style content.}
\centering
\begin{tabular}{lcc}
\toprule
Method & Score & Time \\
\midrule
Placeholder A & 0.91 & 1.4 \\
Placeholder B & 0.88 & 1.7 \\
Placeholder C & 0.84 & 2.0 \\
\bottomrule
\end{tabular}
\end{table}
\section{Conclusion and Discussion}
This page verifies the imported official Springer Nature template path while preserving the six-part manuscript skeleton requested for the library.
\begin{thebibliography}{9}
\bibitem{ref1} A. Smith and B. Lee, ``Template-aware Springer Nature manuscript preparation,'' \textit{Springer Placeholder Journal} (2025).
\bibitem{ref2} C. Garcia and D. Patel, ``Figure and table balance in journal article templates,'' \textit{Scientific Reports} (2026).
\end{thebibliography}
\end{document}
